\section{Benchmark schema and glossary}\label{supp:schema}

The leaderboard is built from a flat record schema. Stating it removes the ambiguity
about what a single number on the benchmark counts. Each method is run per game and writes a
per-game record under its phase's \texttt{out/} tree. The leaderboard aggregates those
records into one row per method by averaging over games. The schema columns of a leaderboard
row are given in Table~\ref{supp:tab:schema}. The aggregation rule is a mean over the
$n_{\text{games}}$ games. When a method emits several output records on one game, those records
are first reduced within that game. The faithfulness, sufficiency, and minimality columns
report the F, S, and M axes defined in \Sectionname\ref{sec:triad}. We do not re-derive those definitions here.

\begin{table}[t]
\centering
\caption{\textbf{Leaderboard row schema.} The columns of one row in
\texttt{compare/out/leaderboard.json} (\texttt{rows[\,]}) and its uncertainty companion
\texttt{leaderboard\_ci.csv}. The metric column names the per-method scoring function from
\Sectionname\ref{supp:protocols}; F, S, and M are the triad axes of \Sectionname\ref{sec:triad}. Aggregation is a mean
over $n_{\text{games}}$, with bootstrap confidence intervals over games (and, where present,
over records).}
\label{supp:tab:schema}
\footnotesize
\begin{tabular}{@{}p{3.0cm}p{8.4cm}@{}}
\toprule
Column & Meaning \\
\midrule
\texttt{phase} & phaseA\_kording, phaseB\_attribution, phaseC\_mechanistic, or ground\_truth \\
\texttt{method} & method identifier (the row key) \\
\texttt{tradition} & intervention, causal, gradient, correlational, \jkey{dim_reduction}, probing, descriptive, oracle \\
\jkey{n_games} & number of games the method was scored on (up to 42 for the scored set of \Sectionname\ref{supp:scoredset}; 1 for the oracle) \\
\jkey{n_records} & number of per-game output records aggregated into the row \\
\texttt{faithfulness} (F) & faithfulness vs.\ the intervention oracle (\Sectionname\ref{sec:oracle}), $[0,1]$, higher better \\
\jkey{faithfulness_ci95} & $95\%$ bootstrap CI half-width over games \\
\jkey{faithfulness_position_regime} & F restricted to the discrete sprite-position/index regime \\
\jkey{faithfulness_content_regime} & F restricted to the continuous content regime \\
\texttt{triad\_S} / \texttt{triad\_M} & sufficiency / minimality axes (\Sectionname\ref{sec:triad}), where applicable \\
\jkey{plausibility_proxy} & documented method-tradition proxy, $[0,1]$ --- not a human-subjects measurement \\
\jkey{is_positive_control} & 1 for the oracle row, 0 otherwise \\
\jkey{metric_names} & the per-method scoring function (\Sectionname\ref{supp:protocols}) \\
\texttt{commits} & the commit(s) at which the per-game records were produced \\
\bottomrule
\end{tabular}
\end{table}

We fix the vocabulary that the schema implies, because the review asked for a single
unambiguous method count. A \emph{method} is one interpretability technique, e.g.\
\texttt{vanilla\_saliency}. A \emph{method row} is its single aggregated line on the
leaderboard. A \emph{per-game record} is one method scored on one game, written under a phase
\texttt{out/} tree. An \emph{output record} is a per-game record for one output regime; a
method may emit several of these on one game. A \emph{family mean} is the average over the
members of a tradition family, e.g.\ causal/intervention. The \emph{oracle positive control}
is the exact $\mathrm{do}(u)$ ground truth of \Sectionname\ref{sec:oracle}, scored as a method to anchor the
ceiling. A \emph{not-applicable method} is one that does not apply to a given output regime;
it is recorded as N/A rather than as a zero. With this vocabulary the count is unambiguous.
The benchmark scores 30 interpretability methods together with the intervention
oracle as a positive control, for 31 leaderboard rows in total
(\texttt{leaderboard.json.n\_methods}~$=31$), aggregating 1773 per-game records over the
42-game scored set (\texttt{n\_per\_game\_records\_aggregated}~$=1773$): 8 methods in
Phase~A, 12 in Phase~B, 10 in Phase~C, and the one oracle row. This count is the canonical one
used in the abstract, the figures, and throughout this Supplement.
